\begin{refsection}
    \renewcommand{\thefigure}{\arabic{figure}}

    \chapterTwoLines{Uma história de vida nos sertões}{o testamento de Francisco Xavier da Costa, Capitania do Rio Grande do Norte (1776)}
    \label{chap:historia-vida-sertoes}
    
    \articleAuthor{Thiago do Nascimento Torres de Paula}
    {Pós-Doutor em Educação pela UFRN (PNPD/CAPES/2017-2018), Doutor em História pela UFPR, servidor da FAPERN (Fundação de Apoio à Pesquisa do Estado do Rio Grande do Norte). Atua como Pesquisador do LEHS/UFRN (Laboratório de Experimentação em História Social da Universidade Federal do Rio Grande do Norte), Docente Colaborador na Pós-Graduação Lato Sensu do IFRN e do Núcleo de Capacitação de Professores da SEEC-RN. Além de autor de inúmeros artigos científicos e capítulos de livros. Lattes ID: 1215.9127.7257.3170. ORCID: 0000-0002-4481-4327. E-mail: thiagotorres2003@yahoo.com.br.}

    \section{Introdução}

    O objetivo  \marginpar{
        \begin{flushleft}
            \tiny \sffamily
            Como referenciar?\\\fullcite{SelfTorresdePaula2021Uma}\mybibexclude{SelfTorresdePaula2021Uma},
            p. \pageref{chap:historia-vida-sertoes}--\pageref{chap:historia-vida-sertoesend},
            \journalPubDate{}
        \end{flushleft}
    } deste trabalho, é apresentar aos pesquisadores em História a transcrição do testamento do senhor Francisco Xavier da Costa, buscando preservar o documento original elaborado na segunda metade do século XVIII, que atualmente encontra-se sob a guarda do arquivo do Instituto Histórico e Geográfico do Rio Grande do Norte.

    Documento normalmente elaborados ao fim da vida, deveria ser escrito, ou ditado por pessoas que obrigatoriamente estivessem em perfeita saúde mental. O estado de consciência do testador deveria ser declarado nas linhas elaboradas, com isso, os testamentos dos séculos XVII, XVIII e XIX assumiram dupla função: permitia o acerto de contas do indivíduo com mundo celestial e humano, mas sobretudo possibilitava a transmissão de bens moveis e imóveis. 

    Poucos foram os testamentos da Capitania do Rio Grande do Norte que chegaram ao tempo presente. O documento que expressa as últimas vontades do Francisco Xavier da Costa (1776), agora transcrito sob as normas do Arquivo Nacional Brasileiro, vem a público com as marcas de sua trajetória, alguns pontos de ilegibilidade que foram causados pela ação de insetos e acidentes com líquidos ao longo do tempo. No entanto, tais lacunas não interferem na compreensão total do texto. 

    Francisco Xavier da Costa era português, como tantos outros reinóis que chegaram ainda jovens na América, contraiu matrimonio na Vila de Goiana, Capitania de Pernambuco, mas finda a vida na condição de viúvo. Tornou-se fazendeiro nos sertões, proprietário de gado, senhor de escravos, dono de casas de morada na Serra do Martins e Vila de Goiana, dono terras e objetos de prata, incluso louças vindas do Oriente distante. 

    Por fim, é consenso entre os historiadores que o testamento elaborado no passado, apresenta-se como fontes importantes e necessárias para compreensão de múltiplas dimensões da vida cotidiana. O testamento em sua estrutura é portador de informações que podem ser examinadas por procedimentos qualitativos e quantitativos \cite{MARCÍLIO1983morte, RODRIGUESAndDILLMANN2013Desejando, SANTOS2013Por}. Em suma, a transcrição em tela, servirá como material para o processo de ensino e aprendizado de jovens pesquisadores, como fonte para investigações no campo da História Sociocultural, História Econômica e História da Educação.  

    \section{Transcrição}

    \vspace{5mm}

    Testamento de Francisco Xavier da Costa (1776).

    \vspace{1em}

    \noindent{}[\textit{Documento com lacuna no início}]

    \vspace{1em}

    \noindent{}[3 linhas ilegíveis] dez mil réis = Deixo para as obras do [ilegível] dois mil réis = Deixo [ilegível] se achar [ilegível] pela [2 linhas ilegível] filho [ilegível] na [linha ilegível] natural da Vila de Pudantes, Bispado de Coimbra ambos já defuntos. Declaro que do casal de meus pais não [ilegível] saí com [ilegível] de [ilegível] vindo para a cidade de Lisboa de lá vim para [ilegível] Brasil com idade de vinte anos pouco mais ou menos. Declaro que se [ilegível] morte de meus pais me tocou alguma coisa o poderá saber meu irmão, o reverendo padre beneficiado Antônio Xavier da Costa que foi a que tomei a sua conta nos particulares da casa por morte de meus pais. Declaro que casado na vila da Goiana, Bispado de Pernambuco, com Maria de Castro, filha legítima de João de Castro França, da sua mulher Maria [ilegível] já defuntos. Declaro que tive da minha mulher dois filhos, uma por nome Mariana, e outro por nome Antônio que este faleceu de idade de quatro anos pouco mais ou menos. Declaro que por morte de minha mulher se fez inventário de seus bens, na qual [ilegível] ram mais filhos suas meações, as quais meações tocou a [ilegível] filho um escravo por nome Pedro crioulo, o qual ficou em meu poder como administração do dito meu filho = Declaro que este escravo por nome Geraldo me fugiu há bastante anos e se este aparecer ou o seu produto pertence a sua [ilegível]  a dita minha filha [ilegível] herdeira do dito seu irmão = Declaro que a dita minha filha Mariana Xavier da Costa [ilegível] carregando a Maria escrava ficando nesta forma inteirada [ilegível] morte de minha [ilegível] escrava do gentio da [ilegível] nome Joana [ilegível] puder da dita minha [ilegível] ficou um filho por nome [ilegível] filha o qual escravo quando fugiu a dita minha filha para se casar contra minha vontade a não quis acompanhar valendo-se de mim para eu fazer com ele e o não querendo eu que buscava senhor para a comprar temendo-se do castigo porque o ameaçava o dito senhor por não fazer acompanhar a senhora quando fugiu o qual escravo eu nunca pus encarregar e fiz com que o dito senhor me passasse um papel de venda por trespasso de uma escrava que a dita minha filha tinha levado, o que fez obrigado, e com temor, dizendo que por minha morte sempre devia haver a si o dito escravo = Declaro que depois de feito este papel de venda temendo-me o dito escravo Fabrício que por minha morte o puxassem a si para o castigarem entrou este a meter respeitos para que eu repasse sua carta o que pereça de muitos rogos repassei de minha letra sem condição alguma por todos me dizerem, que dita escravo me havia acompanhar e servir enquanto fosse vivo e o dito assim a confessava e passando-lhe eu a dita carta lhe não quis entregar, e sim a dei a guardar a quem a pudesse se por minha morte o quisessem cativar = Declaro que falecendo o senhor marido da minha filha e sabendo o dito escravo da dita carta entrou a meter empenhos o que eu a lançava nas notas dizendo me haverá servir e acompanhar enquanto eu fosse vivo, o qual na dita carta em meu poder e passada poucos [ilegível] estando eu no sertão na minha fazenda e fixando ao dito escravo [ilegível] casa da minha residência se ausentou sem causa alguma carregando-me várias coisas de casa no que tem desobedecido e perdido a dita liberdade de não poder valer dela pela desobediência e furto que me fez e não ter o dito meu genro recebido coisa alguma pelo valor do dito escravo faço esta declaração para desencargo de minha consciência e poderem haver [ilegível] si os senhores querendo [ilegível] se o dito escravo [ilegível] e fizesse o furto a minha [ilegível] pelo seu valor pelo dito meu [ilegível] escravo // Declaro que [linha ilegível] a sua [ilegível] Maria Pinto de Araújo que sua testamenteira a que aceitou a [ilegível] lhe se deu sua meação depois de [ilegível] para pagar [ilegível] ficando a dita testamenteira [ilegível] para a dívida como o da meação da testadoura na conta [ilegível] feito na justiça [ilegível] de Porto Alegre dando logo cumprimento [ilegível] disposições testamentárias e consta no livro das rendas [ilegível] as ditas contas. Declaro que depois de dar contas dos bens [ilegível] testamentária da maior parte dele a dita testamenteira D. Maria Pinto de Araújo [ilegível] a nulidade do dito testamento com meus irmãos o qual testamento se houve por nulo os ditos herdeiros habilitados por herdeiros dos bens que existissem, dos quais bens estava de posse para mim e como procurador bastante de meus irmãos assistentes estes em Portugal. Declaro que por morte da dita testamenteira, D. Maria Pinto de Araújo, estavam todos os bens de gado e bestas pró indiviso, e no inventário que se fez para [ilegível] da dita testamenteira D. Maria se partiram estes ao meio entre a dita testadoura D. Maria e os herdeiro do dito coronel Antônio Duarte Teixeira, com os quais bens levaram aos herdeiros do dito coronel Antônio Duarte Teixeira tenho pago várias dívidas que devia a meação do dito coronel Duarte como fosse duzentos e tantos mil réis aos herdeiros de José Gonçalves da Silva como consta da escritura da composição que se foi com os ditos herdeiros, a qual se acha nas atas da justiça da Vila Nova de Porta Alegre, no sertão do Apodi e outras mais dívidas que tenho pago = Declaro  que ainda a meação dos bens do dito [ilegível] Antônio Duarte está devendo quinze mil réis [ilegível] do padre Tomé e estes se não sabe quem sejam a toca a dita dívida entregar-se aos cujos donos ou ao Tribunal dos Ausentes, e quando eu não tinha satisfeito e pago esta dívida meus testamenteiros o averiguaram  a quem se deve pagar e a pagaram logo não tendo o recibo de que tenho pago. Declaro que por morte de meus pais não tenho herdado coisa alguma, e meus testamenteiros averiguados saber de meu irmão o padre beneficiado Antônio Xavier da Costa, morador na Vila de Penela, bispado de Coimbra. Se me tocou alguma coisa da herança e sendo este falecido averiguaram saber de minha irmã Angélica Maria de Jesus, assistente na mesma vila de seus herdeiros, sendo eu o não tinha [ilegível] minha vida e sendo haja de me ter tocado a alguma coisa e averiguara o quanto estarão os meus testamenteiros pelo que diferem os meus irmãos ou os seus herdeiros e achando que haja de me tocar alguma coisa aos ditos meus irmãos ou seus herdeiros da herança do dito meu tio Antônio Duarte Teixeira, de cujos bens está de posse e não tendo a coisa que me possa tocar meus testamenteiros averiguarão o que lhe poderá tocar aos ditos meus irmãos de herança dos bens do dito meu tio Antônio Duarte para lhe fazerem remessa quanto eu o não tenha feito em minha vida, de cujos bens de que estou de posse são os seguintes = Declaro que tomei posse dos bens da herança de uma casa na Serra do Morez, os quais foram avaliados em sessenta mil réis por ser uma armação velha, cujas casas que se não perderam as consertei e nelas tenho gasto perto de duzentos mil réis, sendo que em minha vida as não tenho vendido dela se tocara o valor de sessenta mil réis e o mais velho quem forem vend[ilegível] vira para a morte do meu casal para se fazer partilhas com me [ilegível] terça  = Declaro que tomei posse de uma data de três léguas de terra no Panema chamada a Pedra cumprida, a qual vendi por quinhentos e sessenta mil réis com o qual paguei a dívida que tenho declarado e o mais está em meu poder. Declaro que tomei posse de uma sorte de terras no sítio chamado o Saco que seu valor consta do inventário que se fez dos bens que ficaram do dito meu tio Antônio Duarte Teixeira e o valor da dita terra [ilegível] do mesmo [ilegível] na Justiça da vila de Porta Alegre = Declaro que também estou de posse do gado que tocou ao dito coronel, de sorte que se acha situado em ditas terras do Jacó, o que escapou da seca de setenta e oito = Declaro que deste gado tenho disposto alguns bois que não sei o compito por se não fazer assento nas listas por ignorância ou descuido = Declaro que dos bens de que estou de posse me não poderá coisa alguma da minha parte da herança pelo que poderei ter disposto e me utilizado dele e se venderão os bens que existirem quando eu o não tenha feito em minha vida e do seu porte meus testamenteiros farão remessa aos ditos meus irmão ou seus herdeiros quando eu o não tenha feito em minha vida que há no que cuido para de encargo desencargo de minha consciência = Declaro que fui testamenteiro de Maria Pinto de Araújo e dos seus bens fiz inventário na justiça de Porta Alegre, do qual testamento não cheguei a dar contas por se oporem a validade do dito testamento seus herdeiros, como fosse D. Tereza Tavares de Jesus assistente na cidade da Paraíba e outras mais, os quais levaram sentença pelo preso dos resíduos = Declaro que em verdade da dita sentença fiz entrega de todos os bens a dita Tereza Tavares de Jesus, como consta do termo de entrega e contas aprovada e assinada pela dita D. Teresa Tavares e seu agente e procurador, o tenente general Antônio Xavier Tavares da Silva, julgado este termo por sentença que se acha acostada ao dito inventário = Declaro que por morte da herdeira D. Tereza Tavares de Jesus fez este seu testamento deixando por sua herdeira Tereza Tavares e a seu filho, o tenente general Antônio Xavier Tavares da Silva, por seu testamenteiro, o que melhor constará do dito testamento, e estes estão de posse de seus bens pertencentes a dita herança que teve D. Teresa Tavares de Jesus dos bens de Dona Maria Pinto de Araújo do qual encargo se dá contas no juízo dos Ausentes em Pernambuco e inventário que se fez dos bens da dita D. Maria Pinto de Araújo se adjudicaram bens para este encargo conforme um mulato por nome Miguel e deu contas e a [ilegível] de gado e outros mais bens para pagamento das dívidas que o casal devia que tudo consta do inventário que se foi [ilegível] da vila de Porta Alegre de cujos bens tomou para a dita herdeira D. Tereza Tavares de Jesus para ir satisfazendo ao dito encargo como herdeira que foi da dita D. Maria Pinto de Araújo Declaro que quando foi entregue e dei contas a dita herdeira D. Tereza Tavares estava em dito casal de D. Maria Pinto de Araújo quinze certidões de capelas de Missa pelo encargo de Monica Pinto Coelho, de que se não tinha dado conta delas em dito juízo das quais certidões fiz entrega delas ao tenente general Antônio Xavier Tavares da Silva, como procurador bastante e agente da dita D. Teresa Tavares de Jesus por ser o que tudo fazia naquele casal de quem era procurador para com as ditas certidões dar conta naquele juízo com os mais que se fosse seguindo, quando lhe fossem pedidas. Declaro que tenho noticio que a dita herdeira D. Teresa Tavares de Jesus era herdeira e testamenteira nas [ilegível] em contas deste encargo, mas antes tem disposto herdeira adjudicados para o dito encargo como fosse em forra [ilegível] Miguel por cento e vinte mil réis, preço que que lhe deram em o dito inventário achando quem desse por ele se fosse a praça duzentos mil réis por ser este bom vaqueiro e a gado o tem disposto e o resto vendido ao licenciado levas da Paraíba que [ilegível] vaqueja e ferrar pelo seu vaqueiro na fazenda dos Currais Velhos, onde estava situado = Declaro que faço esta [ilegível] constar e desencargo de minha consciência = Declaro que ando tratando dos particulares das dependências do casal de meu tio o coronel Antônio Duarte Teixeira, casado com D. Maria Pinto de Araújo e chegando da Bahia onde tinha ido tratar de várias demandas e vindo algumas vencidas neste meio tempo, faleceu o dito meu tio [ilegível] e mandando-me a dita sua mulher assinar do sertão onde eram moradores fui para ajustar as minhs contas com a dita D. Maria do que tinha dispendido da minha fazenda [ilegível] a dizer dependências [ilegível] e chegando [ilegível] o dito [ilegível] contas, e fazendo [ilegível] das duas dependências quando [ilegível] de [ilegível] forma que a dita D. Maria dizendo-me que [ilegível] desamparava pedindo-me que o dito casal [ilegível] que não [ilegível] para me pagar ao que me [ilegível] e por empenhos e rigor assistisse eu com ela para lhe trazer as suas dependências como há tantos anos o tinha feito em vida de seu marido, o meu tio o coronel Antônio Duarte e que se eu o desamparava pelo seus anos tudo [ilegível] sem ela lhe poder dar remédio dizendo-me que aceitasse como pagamento em gado que ela que me dava um pedaço de terra que tinha no sítio do Saco para eu ter onde situar meus gados, ao que me sujeitei, ainda que com pouca conveniência e por a dita D. Maria Pinto não ter os [ilegível] forçados que se me dessem sua fazenda me fez escritura de doação remuneratória da dita terra com a obrigação de eu estar na sua companhia enquanto ela fosse viva e tratar das suas demandas e dependências de seu casal, assim como o tinha feito havia tantos anos em vida de seu marido, o meu tio o coronel Antônio a dita escritura de terras para da dita terra, e nela entrei a situar os gados de que me fazia pagamento para o que me devia por ordem que mandou o capitão Antônio da Rocha de Carvalho para me fazer o dito pagamento = Declaro que fiquei na dita casa tratando de todas as dependências do dito casal como procurador bastante que me fez a dita Maria, fazendo por ordem sua composições nas ditas demandas que haviam tantos anos que corriam vendo do sertão a esta praça a compor várias demandas como fosse uma que corria havia uns vinte tantos anos com o testamenteiro de seu pai Miguel Pinto, e outras com o tenente José Rodrigues Chaves em Goiana redundando-lhe muitas composições, muitas conveniências e sossego que a não ser assim não teria nenhuma mortalha para se amortalhar e pagando dívidas que o casal devia fazendo tudo o que era da minha obrigação, enquanto a dita D. Maria foi viva = Declaro que falecendo a dita D. Maria Pinto com o seu solene testamento [ilegível] D. Teresa Tavares de Jesus a validade do dito testamento e outros mais pela lei novíssima pela dita D. Maria instituir a sua alma por herdeira do remanescente da sua fazenda, da qual tiveram sentença havendo-se o dito testamento por nulo e que se entregassem todos os bens aos herdeiros habilitador = Declaro que com esta sentença fiz entrega de todos os bens a dita herdeira D. Teresa Tavares de Jesus que tudo consta do inventário dos ditos bens = Declaro que nas contas que dei e entrega que fiz dos bens das contas constava estarem o dito casal devendo trinta e nove mil quinhentos e setenta e quatro réis que logo o procurador e agente, o tenente general Antônio Xavier Tavares da Silva me pediu passasse eu recibo  ditas contas o que fiz sem que nada tivesse recebido = Declaro que ajustando as ditas contas com a dita D. Maria por seu procurador e agente, o tenente general Antônio Xavier Tavares da Silva me saio o dito com uma dívida de gastos que tinha feito a que eu devia pagar dez mil e tantos réis de um chapéu que me tinha comprado um galão que empostava cinco mil e tantos réis, e quatro mil réis por conta de oito mil réis que tinha pago por conta das dependências do casal que eu era obrigado a pagar a meta de oito mil réis que pagavam ao fojes letrado no que eu a nada pus dúvida = Declaro que depois de justas as contas e feita entrega de tudo, a dita herdeira e a seu procurador e agente, e feita entrega de todas as certidões das capelas de missas do dito encargo de Monica Pinto Coelho, me saiu o dito tenente general dizendo-me que a dita D. Maria Pinto de Araújo não podia fazer a tal doação de terra, porque a terra estava avaliada em duzentos mil réis e que a doara não podia dar cinco até tanto dizendo-me me punha demanda como procurador da dita D. Teresa e por eu evitar estar dei ao dito tenente Antônio Xavier como procurador e agente sessenta em dinheiro, digo, em dinheiro sessenta um mil e tantos réis com o dito que se me devia do soldo das contas que importou trinta e nove mil quinhentos e setenta e quatro réis, que uma e outra conta importava em oitenta e tantos mil réis que recebeu me passasse recibo de recibo de nada do que recebeu, querendo inda que eu lhe desse mais da conta que não podia dever sem que a dita herdeira D. Maria disto de nada soubera = Declaro que depois de morta D. Teresa deixou por seu testamento ao dito Antônio Xavier e sua mãe por herdeira entrou o dito Antônio Xavier Tavares e persuadiu-me a que lhe desse mais dinheiro por conta da dita terra doada. E aconselhando-me eu com letrados me disseram que nada lhe devia dar, porque a dita doadora quando doou a dita terra não tinha herdeiros forçados e ao depois os tios necessários pela lei novíssima e como não eram herdeiros forçados podia em sua vida doar tudo que possuía e fazer dele o que quisesse e quando o dito tenente Antônio Xavier queira demover alguma dúvida aos meus herdeiros estes a chamaram a juízo para jurar se recebeu as ditas parcelas e certidões de missas e quando queira negar tem testemunhas que sabem deste dinheiro que lhe dei como seja o capitão-mor Francisco Arruda e João Martins Bragança que este assistis na ocasião na casa da viúva de Felipe Gomes declaro na cidade da Paraíba, quando dei as contas e fiz entrega dos ditos bens e dinheiro = Declaro que me deve José Antônio da Silva, escrivão que foi da povoação do Assú, dinheiro de empréstimo que me mandou pedir do Recife quando foi rematar o seu ofício por uma vez que mandou pedir trinta mil réis e lhe mandei cinco meias doblas e logo de outras vez me mandou pedir cinquenta mil réis lhe mandei oitomil deles que por não ser dinheiro trocado na ocasião lhe mandei demais em dita [ilegível] que tudo consta dar cartas eu tinha em meu poder. Declaro que vindo Antônio José da Silva com ofício rematado por esta vila, buscando a minha casa com o protesto dizendo-me que no sertão o tornaria a dar como também pagar-me o que me devia de dinheiro de empréstimo. Declaro que não haviam que o dito Antônio José veio do [ilegível]  eram [ilegível] nova em que ia dar em pagamento do que me devia em preço de cinquenta mil réis e pelo mesmo preço a ofereceu a algumas pessoas pelos mesmos cinquenta mil réis para mudar a conta do me devia e não achando quem a quisesse pelo preço me deixou em dizendo-me que se eu a pudesse levar quando fosse para o sertão a fizesse e na falta a deixasse ficar na casa da viúva Maria de Perez = Declaro que indo eu para o sertão com uma filha viúva que tinha vindo buscar para minha companhia e levando a outros escravos pequenos levei a dita negrinha e logo que cheguei na minha casa fiz aviso ao dito Antônio José da Silva no Assú, para que viesse buscar a dita negrinha e me trouxesse o que me devia, o que nunca fez = Declaro que passados anos veio o dito Antônio José da Silva da povoação do Assú, aonde eu assistia na minha casa na Serra de Martins, e falando-lhe eu no que me devia me deu a dita negrinha em preço de sessenta mil réis, dizendo-me que o resto que ficava devendo me mandava logo do Assú, o qual nunca mais vi, por anda criminoso = Declaro que adoecendo no sertão e vindo-me a curar em Goiana tive notícias que o dito Antônio José da Silva se tinha ausentado das partes do Assú, deixando a muitos fintados, e que dela buscara a Serra de Martins, e que ocultamente carregara a dita negrinha levando-a para as partes dos Piauís, para a vender fazendo furto e deixando de me pagar o que me devia que de tudo sabe minha herdeira, e poderá cobrar a dita dívida sendo que o dito apareça em algum tempo, e eu o não tenha cobrado = Declaro que tenho contas com o defunto Bento José de Bessa, a quem emprestava alguns dinheiros para negociar seus gados e dando-lhe os meus gados de frete para Gregório José e na ausência a entregar ao capitão-mor Álvares para me mandarem dispor, dando ao dito ordem para me comprar algumas encomendas carecia, estando sempre pelas contas que o dito me dava = Declaro que vindo eu para Pernambuco em ocasião que o dito Bento vinha também com uma boiada que tinha negociado com o sargento-mor Antônio Pereira do Itapuã, ali me pediu o dito Bento Jose [ilegível] cem mil réis para deixar de sinal de outra boiada que lhe emprestei, presenciando isto André de Souza que vinha comigo na ocasião a fazer seu negócio em Pernambuco = Declaro que vindo eu para esta vila, e me arranchei nela alugando eu umas casas para assistir nelas e chegando o dito Bento José de Bessa com a sua boiada, dispôs nesta vila arranchando-se na minha casa pela amizade que comigo tinha = Declaro que dispondo o dito Bento a dita boiada se foi para Pernambuco a cobrar o produto dela, e fazer suas disposições e pagamentos passados vários dias chegou da praça com seus negócios e fazendas suas e alheias arranchando-se na minha casa e queixando-se de que não tinha feito negócio e conveniência com a dita boiada me cometeu o negócio para comprar farinhas para levar para o Aracati no que eu convim, sendo sócio com o dito no negócio das farinhas, e entrando eu e o dito a comprar farinhas e panos de algodão e sacos para ensacar e a embarca-la, sucedeu ao dito Bento José de Bessa cair doente de bexigas na minha casa, onde da dita moléstia deu alma a Deus fazendo seu testamento e aprovando = Declaro que logo que caiu doente das ditas bexigas fez seu testamento e nele declarou com quanto tinha entrado e eu ajustando todas as contas que tinha tido comigo, pelas quais estava por estarem por sua letra e eu não vi clareza alguma mais do que ao que ele dava, cujas estão em meu poder = Declaro que estando parte das farinhas embarcadas falei a Manuel de Figueiredo as ir dispor dando-lhe um escravo meu para o acompanhar, ficando o dito Bento com a dita moléstia, da qual deu contas a Deus = Declaro que falecendo o dito com o seu testamento me nomeou nele primeiro testamenteiro, o qual não quis aceitar, mas do que para lhe fazer as disposições do seu enterro, e o mandei registrar o dito testamento e remeti a seus herdeiros por um seu tio por nome o capitão Francisco Saraiva de Moura que também na ocasião tinha vindo com boiada, e sendo entregou o dito testamento e os ditos herdeiros este logo buscaram o Aracati para onde tinham ido as ditas farinhas e disporem-se e com o dito testamento atemorizaram ao homem que eu tinha manda dispô-las e tudo entregou ficando-se estes com tudo sonegando o dito testamento que deles não deram contas, ficando eu tão somente em meu poder com as clarezas das contas que o dito tinha comigo feito pela sua letra e tudo o mais que tinha em meu poder pertencente ao dito defunto Bento fiz entregar como consta dos recibos que se acham em meu poder e certidões dos gastos que fiz no enterro que averiguarão esta conta os meu herdeiros e testamenteiros para se saber o que se me deve = Declaro que os bens que possuo foram adquiridos e são os seguintes = Declaro que possuo um escravo do gentio da Angola por nome João, e dois moleques Angolas um por nome Antônio, e outro por nome Francisco, uma cabrinha por nome Ana que houve por rematação do casal de José Vieira Lima, uma mulata por nome Maria da Luz que houve por rematação dos bens de José Vieira Lima, para auto de arrematação junto ao inventário que o fez de seus bens no cartório do capitão João Figueira no juízo ordinário da vila de Goiana, e outro mais bens que consta da mesma arrematação = Declaro que possuo um mulatinho por nome Serafim houve por compra que dele me fez D. Teresa Tavares de Jesus, como Herdeira dos bens de D. Maria Pinto de Araújo que consta de papel de venda que me fez a dita D. Teresa que se acha em meu poder, cujo mulatinho me fugiu em Goiana, e dele não tive mais notícia que próximo o ter sonegado sua mãe ou parentes, por estes ter ficado forro a causa morte em dito testamento, visto o anulara a dita herdeira e cativara o dito mulatinho que a comprou por este ser meu afilhado e se aparecer em algum tempo vira para o meu casal para dele se fazer partilha u de seu produto = Declaro que possuo uma escrava do gentio da Angola por nome Rosa que está em meu poder empenhada que consta do papel de empenho que tenho em meu poder  = Declaro que o ouro que possuo é o seguinte: um cordão grosso com o peso dezesseis oitavas e meia e dez graus, e com o dito cordão um a relicário de ouro com o peso de nove oitavas com um santinho dentro = Declaro que possuo um par de botões lavrados de punhos de camisa com o peso os dois pares duas oitavas  = Declaro que a prata que possuo é um espadim de uso lavrado com bocal e pontura  e buldziar e fivelas e argolas que o seu peso constará = Declaro que possuo um xafarate om bocal e ponteira e sua fivela grande que o seu peso constará = Declaro que possuo uma caldeira de prata com sua corrente já usada com peso de oitenta e sete oitavas = Declaro que possuo umas esporas de prata de uso com suas fivelas com o peso de sessenta e duas oitavas = Declaro que possuo umas fivelas maltesas de prata com xarnelhas de ferro com o peso de trinta oitavas = Declaro que tenho dois pares de fivelas de ligas de calção de prata com o peso ambos os dois para de dezesseis oitavas dos quais trago em meu poder = Declaro que possuo uma bolseta de tabaco de prata que o seu peso constará = Declaro que possuo uma fivela de pescoço de gravata de prata com o peso de oito oitavas e meia = Declaro que possuo dois pares de fivela de ligas de prata de ata meia com o peso ambos os pares de seis oitavas e meia= Declaro que possuo mais dois pares de fivelas de sapatos com suas xanilhas de prata, um par feito no reino polidas e outro para se ser polidas que deixou na minha casa na Serra de Martins que o seu peso o dirá = Declaro que possuo mais um par de fivelas de ligas de calças de prata que deixei na casa na Serra, e assim mais uma fivela de gravata que seu peso o dirá = Declaro que possuo meia dúzia de colheres de prata vasadas, e seis garfos de quatro dentes com o peso as colheres de cento e dezesseis oitavas, os garfos com o peso de oitenta e oito oitavas = Declaro que possuo duas facas com cabos de prata com o peso a prata de quarenta e quatro oitavas e meia = Declaro que possuo sete colheres de prata de chapa e sete garfos vasados com três dentes cada um, as quais colheres e garfos deixei três colheres e três garfos na Serra, e trago quatro em meu poder que o seu peso dirá = Declaro que possuo um relógio de pixibeque que me custou dezesseis mil réis = Declaro que o cobre que possuo é uma bacia com dezenove ou vinte libras nova de tomar banhos que o seu peso o dirá = Declaro que possuo dois tachos de cobre usados grandes e um meiam, e dois pequenos de libra nova que todos o seu peso o dirá = Declaro que possuo um caldeirão de aguentar o que com sua tampa já usado que o seu peso o dirá = Declaro que possuo uma xiculateira já usada de aguentar e que para chá dois valos de cobre, um pequeno e outro grande que o seu peso dirá. Declaro que possuo duas bacias de urinar, uma de cobre e outra de latão nova que o seu peso o dirá = Declaro que possuo um candeeiro de latão em bom uso que o seu valor o dirão = Declaro que possuo dois pratos de estanho rasos em bom uso que o seu valor o dirá = Declaro que possuo uma bacia e jarro de água ambas de estanho que é de bom uso que o seu valor dirá = Declaro que possuo um pixel de estanho de fazer chá que o seu valor o dirá = Declaro que possuo um inxão de sopeiras ataubadas de cinco [ilegível] que o seu preço o dirá = Declaro que possuo quatro pratos [ilegível]pratos finos grandes, e dois maiores, de duas dúzias de pratos de guardanapos finos, entrando nestes cinco dourados duas tiplas grandes douradas, quatro tiplas pequenas douradas, umas galhetas de azeite, oito xícaras de casa com seus pires e duas xícaras com seus píeres de tomar chá, tudo da Índia e um bule da Índia de fazer chá que tudo os seus valores o dirá = Declaro que possuo um almofariz de bronze que o seu valor o dirá = Declaro que possuo dois copos de beber água de cobre cada um, o tenho na fazenda e o outro nas  casas da Serra = Declaro que toda esta louça nomeada e mais trastes miúdos os tenho nas minhas casas na Serra, e outros mais trastes de que não tenho lembrança = Declaro que possuo na Vila de Goiana na rua do Rosário dos prestos uma morada de casas de duas salas que pagão foro a Nossa Senhora do Rosário da Matriz onde eu estou assistindo e dentro delas possuo um lanceiro, duas torneiras, dias mesas, uma de duas gavetas , com suas fechaduras e outra de uma gaveta sem fechadura, ambas torneadas = Declaro que possuo em ditas casas um armário partido ao meio, com suas fechaduras = Declaro que possuo mais na dita casa uma caixa grande com duas gavetas e suas fechaduras = Declaro que possuo mais duas caixas grandes, as quais estão no Amparo, em casa de Manuel Figueiredo = Declaro que possuo um oratório pintado de preto, um bocal de duas fechaduras de Mascavia já usado, um estrado, um leito de Jacarandá torneado, um contador = Declaro que tenho emprestado outro contador ao capitão-mor Antônio de Albuquerque em Goiana Grande, que se mandará buscar, e assim mais possuo várias cadeiras velhas que a tudo se dará seu valor quando for tempo = Declaro que possuo outra morada de casas de duas salas que pagam foro a Nossa Senhora do Rosário da Matriz, e ambas as moradas pagam foro a tostão por braça, as quais são citas no beco do Pavão = Declaro que possuo na Serra do Martins uma morada de casas que pagam foro a Nossa Senhora da Conceição da mesma Serra das quais casas se hão de tirar sessenta mil réis, valor que se lhe deu no inventário de meu tio, o coronel Antônio Duarte Teixeira, e vendidas que sejam se hão de remeter as ditas sessenta mil réis a meus irmãos em Portugal, na vila de Penela, Bispado de Coimbra, e o mais que renderem as ditas casas me pertence das benfeitorias que lhe tenho feito, e gastos que gastei nas benfeitorias perto de duzentos mil réis = Declaro que os bens que possuo em ditas casas na Serra são dois machados, uma foice nova, cinco enxadas em bom uso, uma torneira, um ante banco de encosto, outro sem encosto, duas camas de frade cobertas de couro, uma banca, uma mesa com duas gavetas com suas fechaduras um ferro de cova = Declaro que em ditas casas possuo duas frasqueiras, uma com dose frascos, outra com dez frescos, ambas esticadas, dois garrafões encourados = Declaro que possuo três espingardas, duas que tenho nas casas da Serra e uma que trago comigo = Declaro que tenho  na dita Serra uma catana de ferro, e possuo outra frasqueira de oito frascos que trago comigo = Declaro que possuo nas casas da dita Serra uma verde pintada nova, e três brancas pintadas = Declaro que possuo nas ditas casas duas garrafas de vidro branco oitavadas com suas carapuças, e quatro copos grades de beber água, dois de pé de calis de beber vinho e assim mais possuo em ditas casas uma dúzia de pratos de guardanapos brancos dos de Lisboa e vários pratos de estantes ou tigelas = Declaro que os meus herdeiros declararam os demais trastes que tenho de portas e dentro e assim mais possuo nas ditas casas uma peia de ferros de pear cavalos que me custou quatro patacas e uns caixões de guardar cabeleira = Declaro que possuo uma sorte de terras de criar gados no sítio do Saco chamado Currais Velhos que consta da escritura que dele se me fiz que a tenho em meu poder, passada nas notas do Juiz ordinário da Vila de Porta Alegre = Declaro que o dito sítio o tenho sitiado com gados vacuns que o seu número se saberá depois de avaliar chegado a meiram para se saber o que me toca da terça = Declaro que no dito sítio tenho cinco magotes de bestas que o criados poderá saber o número como também os cavalos que há de fábrica e poldros de negócio = Declaro que possuo em dita fazenda um alabanca e um alvião e cinco machado, uma serra e dois ferros de co[ilegível]as, um martelo, um [ilegível]opulo grande e dois pequenos e um [ilegível] ferros [ilegível] dará [ilegível] grandes. Declaro que na dita fazenda tenho dois [ilegível] com suas estribarias de ferro e celas bastardas que tudo que houver declarará o dito criador que estiver em dita fazenda = Declaro na dita fazenda tenho dois [ilegível]lheiros e meio ou três de telha que mandei fazer para umas casas em dita fazenda = Declaro que possuo três celas [ilegível] e a casa da Serra, as quais casas possuo também um lanceiro de ombros de guardar vestidos e uma cabeça de cabeleira com o seu par, a cujos bens se dará valor por minha morte = Declaro que falecendo José Vieira Lima casado com a minha herdeira lhe ficará tantas dívidas que os seus bens não chegarão para pagamento nem das pestificadas que foram antes raptadas e sendo rematados todos os bens para pagamento das ditas dívidas eu os remi fazendo-se-me três por [ilegível] bens rematados, ficando eu obrigado a pagar as dívidas nomeadas, os quais paguei e só não paguei a dívida pestificada de Antônio Batista Filho, por este ter sido herdeiro de meu enteado o capitão Bento Jose de Bessa, com o qual eu tinha contas e me é devedor como consta das contas que tenho em meu poder, as quais averiguarão os meus testamenteiros quando eu o não tinha feito em minha vida, e na que se me dever descontarão a sua dívida e o resto que contar estar-se me devendo virá para o meu casal. Declaro que os bens que aqui nomeio sendo que neles haja diminuição o declararei em codicilho. Declaro que os bens que me tocarem na minha terça meus testamenteiros os disporão em praça ou fora dele em forma que sejam reputados e com menos gastos= Declaro que não deixo declarado no meu testamento alguma pataca que poderei ter por delas compor para meu sustento e alimentar-me em alguma doença que poderei ter = Declaro que sendo por minha morte não tenha a quem haja de tocar a minha fazenda os meus testamenteiros logo toda a minha fazenda para as disposições de pagar a minha dívida e dar posições do meu testamento e pagar as minhas dívidas e legados dos bens que se cobrarem darão a cada [ilegível] que tenho em Portugal filhas de minha Angélica casa com João das Neses de  Oliveira cem mil réis cada uma que as não declaro por não saber quanto são excesso seus nomes o que fará o meu testamenteiro saber para repartir o dito legado, sendo que no meu falecimento eu não tenha herdeiros, ao resto que sobrar se partirá ambos em três partes, uma parte para órfãos pobres, e a outra parte para as recolhidas do recolhimento da validade da Vila de Goiana e outra parte para o sustento dos pobre que se recolhem a hospital da Santa Casa da Misericórdia de Goiana = Declaro que por morte de José Vieira Lima, casado com minha herdeira lhe ficou muita dívida, e pela pobreza em que ficou a dita minha herdeira a vim buscar do sertão e a levei da Vila de Goiana para a minha companhia, a qual por lhe querer dar estado, e ficar amparada por minha morte lhe comprei várias roupas de custo, e ouro que poderia gastar perto de trezentos mil réis = Declaro que estes gastos que fiz de que estava de posse a dita minha filha, e eu lhe não poder dar o estado, este gasto que fiz deve vir com ele por minha morte ao inventário por se lhe achar valor e se o terem a sua meação = Declaro que muitos bens que possuo os não declaro neste meu testamento por me não lembrarem dos quais saberá muito bem a minha herdeira para se declararem no inventário que se fizer deles querendo = Declaro que os bens que ficaram de meu tio, o coronel Antônio Duarte Teixeira de quem sou herdeiro, junto com meus irmãos, inda tem gados espalhados para partir entre os herdeiros do dito meu tio, e os herdeiros de D. Maria Pinto de Araújo, e outros bens que adquiri, como seja uma data de terras na Serra do Apodi, de fronte de São Lourenço em que nos toca a metade do dito sítio, e vou a metade aos herdeiros de D. Maria Pinto de Araújo, e as datas deste sitio confirmadas por El Rei uma que tem os herdeiros da dita D. Maria, a outra que eu tenho que está em poder do capitão José Vieira de Melo de quem poderá haver os meus herdeiros ou testamenteiros e assim tem mais a herança de uma demanda que corre na ouvidoria da Paraíba com o sargento-mor Inácio da Rocha de Carvalho com  testamenteiro dos bens do defunto Francisco de Araújo cuja demanda devem vir dos meus herdeiros e testamenteiro junto com os herdeiros de D. Maria fazer andar para diante por ser nas interferidas em igual parte na metade dos ditos bens e a dita demanda se acha por apelação ao juízo da ouvidoria da cidade da Paraíba. Declaro que me deve por obrigação Manuel Ferreira de Mendonça trezentos e vinte e seis mil duzentos e noventa e cinco réis, dinheiro de empréstimo que quando me não tinha pago ante meu falecimento se há de achar a obrigação em como puder = Declaro que se me deve por obrigação quatro mil réis, cuja obrigação dei ao capitão José Ferreira de Melo, morador no Assú, cuja obrigação ou o produto dela será o dito obrigado a entregar = Declaro que me deve por obrigação Gonçalo de tal sete mil réis ou o que constar da mesma obrigação, o qual se ausentou da Serra do Martins para os Inhamúns, cuja obrigação dei a João Martins Bragança para a cobrar, que será obrigado a entregar a mesma obrigação ou o produto dela = Declaro que me deve por obrigação o reverendo padre João de Paiva cento e dois mil réis e outra mais adições que lhe comprei que constam do meu assento, e assim mais me deve cento e trinta e dois mil réis que lhe entregou o meu vagueiro e desta quantia se deve abater por mim tem pago do escrivão da vila do Porta Alegre e o dinheiro que deve a minha herdeira de que lhe passou recibo que se estará pela conta que der o dito senhor reverendo Padre quando eu não tinha justo contas em minha vida com o dito senhor = Declaro que me deve Felipe Diógenes dois mil réis que lhe emprestei quando foi no Aracati = Declaro que me deve o meu vaqueiro Graciano Pereira resto de dinheiro de empréstimo seis mil e tantos réis ou o que ele na verdade disser = Declaro que sendo obrigada a dita minha herdeira a entrar com o gasto que foi para lhe dar estado e segundo minha lembrança lhe comprei um manto, umas roupas inteiras de cetim preto, um peito de cetim aguluado de renda de ouro, uma saia de seda amarela de toda conta, e vario de ouro, umas braças de colares que ela quiser declarar para desencargo de sua consciência que a minha mente é não lhe tirar do que lhe toca na sua herança, e nem defraudar o que me tocar a minha terça para as minhas disposições = Declaro que revogo se hei por revogado outro qualquer testamento que apareça depois do meu falecimento, fora deste escrito por letra de Manuel Antunes Henriques, por eu lhe ter o ditado, e depois só quero que este valha, e se lhe dê todo o cumprimento e peço e rogo as justiças de Sua Majestade Fidelíssima, tanto seculares, como eclesiásticas, a quem ele possa pertencer lhe deem a sua devida execução por ser esta a minha última vontade para poder ter este meu testamento o seu devido e inteiro cumprimento, torno a pedir e rogar como no principio deste meu testamento pedi e roguei aos senhores padre Leandro Tavares de Góis e ao capitão-mor Gregório José da Silva Coutinho e a Manuel de Figueiredo e ao capitão-mor Tomás Gonçalves dos Santos e o capitão Henriques da Silva Pontes, e ao capitão Antônio Luís do Amaral, e na Serra do Martins  ao reverendo padre João de Paiva e ao alferes Manuel dos Santos Franco que este por esquecimento do escrivão se não declarou no principio deste meu testamento, por me fazerem mercê e serviço a Deus queiram aceitar ser meus testamenteiros para o que os constituo e nomeio por meus tutores, a procuradores com livre e geral administração do meus bens para que os arrecadem, procurem, zelem e feitoriem, tomando e vendendo dos meus bens os que necessário for para meu enterramento e cumprimento de minhas disposições testamentárias, e sendo os queiram pôr em praça o faram nos lugares que lhe parecer o poderem ter maior valor para o que  em tudo os afianço e abono, e hei por afiançados e abonados, tanto quanto em direito posso e lhe concedo todos meus poderes para o referido fim e para o mais que necessário for para clareza dela e de tudo o mais que aqui declarado tenho, que sendo faleça na vida de que minha, e meu corpo será levado para se sepultar [ilegível] e se lhe faça ofício de corpo presente, como dito tenho meu testamenteiros darão a cada pobre que acompanharem [ilegível] como ofício oitenta réis a cada um pobre e ser a todos sacerdotes que se tirem ao dito ofício e homens com postos = Declaro que me acompanharão todas as confrarias e se lhe pagará pelo costume e a mais pompa do meu enterro será ao arbítrio de meu testamenteiro que tudo se lhe deverá em conta nas que der = Declaro que pagas as minhas disposições legados, sendo que sobre fazenda da minha terra, oq eu sobrar se partirá em três montes, e uma parte se dará para aumento do hospital dos pobre da Santa Casa da Misericórdia e outra parte que se repartir para pobres órfãos honrador e brancos ao arbítrio do meu testamenteiro e a outra parte para as obras dos pretos do Rosário da dita de Goiana ou para os recolhidos da soledade da Vila de Goiana para aumento do seu recolhimento, ficando esta esmola a arbítrio do meu testamenteiro para a dar a quem tiver mais necessidade. Declaro que no ano de 60 de outubro, segundo minha lembrança, estando eu no sertão administrando a fazenda de São Lourenço das varges do Apodi de meu tio, o coronel Antônio Duarte Teixeira, onde a tinha situado a venda com as suas bestas, e uma parte de gado vacum do ferro da [ilegível] pelos serviços que eu lhe tinha feito me ordenou que ferrasse eu aparição das bezerras fêmeas daquele ano para minha filha Mariana Xavier, o que assim fiz pondo-lhe nova marca que segundo minha lembrança se ferraram perto de quarenta bezerros que a conta há de constar do livro das partilhas o número das fêmeas tirado as sortes do vaqueiro. Dízimo cujo [ilegível] está em meu poder já veio rubricado pelo capitão Manuel Carvalho de Araújo = Declaro que neste mesmo ano no fim veio uma seca tão grande que morreram quase todas as mais crias tanto naquela fazenda, como nas mais vizinhas = Declaro que por morrer quase todo o gado que tinha situado o dito meu tio naquele sítio de São Lourenço, largou de pagar renda e alguma cabeça de gado que aparecia a mandou botar no sítio das [ilegível] onde tinha outra fazenda povoada, [ilegível] foram umas duas ou três [ilegível] minha filha e tanto [ilegível] gado na fazenda dos currais o mandei sitiar na fazenda dos Currais Velhos, e nessa conduta foi uma cabeça pertencente a dita minha filha, lhe mandei botar o meu ferro que por não ser conhecido o que a dita veio tinha e não levar descaminho que faço esta declaração para constar e desencargo de minha consciência = Declaro que concedo ao testamenteiro que assistir o meu testamento para dela dar conta fora do tempo da lei dois anos mais, e ao dito testamenteiro que aceitar este meu testamento fora da sua vintena lhe deixo delegado dos bens da minha terça quarenta mil réis em dinheiro pelo trabalho que poderá ter em ir assistir ao sertão as avaliações naquele sertão e não chegando os bens da minha terça para o dito legado haverá repartição nos mais bens e não neste = E para quando esta é a minha última vontade do modo que tinha dito, pedi e roguei a Manuel Antunes Henriques este por mim escrevesse e depois de escrito pelo achar certo como o havia ditado me assinei dos meus sinais costumados como também o fez o escrivão que o escreveu ao meu rogo Vila de Goiana de junho de 28 de 1776 

    \vspace{1em}

    \noindent{}Com o testamento que escrevi\tabto{45ex}Francisco Xavier da Costa
    
    \noindent{}Manuel Antunes Henriques

    \vspace{10mm}

    \noindent{}Aprovação de testamento

    \vspace{1em}

    \noindent{}Saibam quantos este público instrumento de aprovação de testamento ou como em direito para sua verdade de melhor no meu lugar, haja eu dizer-me possa virem sendo no ano no nascimento de Nosso Senhor Jesus Cristo de mil setecentos e setenta e seis, aos vinte e quatro dias do mês de outubro, nesta vila de Nossa Senhora de Goiana capitania de Itamaracá, em casas de morada do capitão Francisco Xavier da Costa, donde eu tabelião adiante nomeado e assinado vim aí para como o dito Francisco Xavier da Costa, pessoa que reconheço pela própria de que se achava de pé, em seu perfeito juízo do que dou minha fé que ele me foi entregue o seu testamento, feito em dezesseis, digo, feito em oito folhas de papel com duas meias folhas cheias por ante as laudas de cada uma menos a última que [ilegível] a última lauda desta por acabar [ilegível] testamento a fé a  [ilegível] e meia e pelas quais [ilegível] testador e os mais testamenteiros do mesmo [ilegível] Manuel Antônio [ilegível] dos [ilegível] o aprovasse [ilegível] como [ilegível] a qual testamento [ilegível] da sua escrita [ilegível] vício, borrão nem entrelinha, mas que [ilegível] na última [ilegível] da dita folha em principio dela um borrão mais pare que [ilegível] e nada mais pude que, e por me dizer o dito testador em presentes das testemunhas que mandara escrever este testamento por Manuel Antunes Henriques e depois de lido e lera o achara a seu contento por estar conforme ele mesmo o havia ditado de sua [ilegível] o havia por aprovado, eu tabelião o aprovei e hei por aprovado, tanto quanto devo e posso por bem do ofício que ocupo e sou obrigado sendo [ilegível] conforme as reais ordens de Sua Majestade Fidelíssima [ilegível] fiz este instrumento de aprovação, em que assinou ele dito testador, com as testemunhas seguintes o alferes João Joaquim de Figueiredo, o licenciado Luís Cavalcante Cavalcante de Albuquerque, o capitão Cirilo Rodrigues de Góis, Manuel Ferrei de Mendonça, José Sebastião Lopes Lourenço [ilegível] Manuel Antônio Teixeira de Carvalho, moradoras nesta vila pessoas que reconheço por próprios do que dou minha fé, Eu, João Teixeira Monteiro Cavalcante, digo, Cavalcante, tabelião que escrevi.

    \vspace{1em}

    \noindent{}\hspace{45ex}Em fé de verdade.

    \vspace{2em}

    \noindent{}\hspace{45ex}\begin{tikzpicture}
        \draw[black, ultra thick] (0,0) -- (.5,0);
        \draw[black, ultra thick] (.25,.25) -- (.25,-.75);
    \end{tikzpicture}

    \vspace{2em}

    \noindent{}\hspace{45ex}João Teixeira Monteiro

    \noindent{}Francisco Xavier da Costa

    \noindent{}Luís Carvalho de Albuquerque\tabto{45ex}Cirilo Rodrigues de Góis

    \noindent{}José Sebastião Lopes\tabto{45ex}Lourenço
    
    \noindent{}Antônio Teixeira de Cavalcante

    \noindent{}João Joaquim de Figueiredo

    \noindent{}Manuel Antônio Teixeira de Carvalho

    \vspace{10mm}

    \noindent{}Termo de abertura

    \vspace{1em}

    \noindent{}Aos quinze dias do mês de março de mil setecentos e oitenta e oito anos, nesta Serra do Martins, abri este testamento do defunto Francisco Xavier da Costa, por pertencer a alternativa e tentativa e para o testamenteiro, o sargento-mor João Dias Martins dar logo cumprimento aos legados para o que entreguei. Serra do Martins, dia e era supra.

    \vspace{1em}

    \noindent{}\hspace{45ex}Francisco Benício de Carvalho

    \noindent{}\hspace{45ex}Vigário da Vara do [ilegível]

    \vspace{1em}

    \noindent{}Aceito este testamento, como testamenteiro que sou, como consta do codicilho, para dar cumprimento aos legados, percebendoeu vintena, e o mais que neste testamento declarado fica ao testamenteiro. Serra do Martins, 17 de março de 1788

    \vspace{1em}

    \noindent{}\hspace{30ex}João Dias Martins 

    \vspace{1em}

    \noindent{}Testamento e última vontade do coronel tenente Francisco Xavier da Costa, aprovado por mim tabelião a qual assina-se e cuzido, lacrado com três pingos de de lacar do reino, zelos azul com outros tantos pingos de lacar vermelho por banda. Vila de Goiana, agosto, 24 de 1788.

    \noindent{Tabelião}\tabto{45ex}João Siqueira Sebastião


    \section{Fonte}

    Fundo Documental do Instituto Histórico e Geográfico do Rio Grande do Norte. Livro de Notas. Testamentos Nº cx. Ant. 113. Nº. Cx atual: Ano: 1776. s/nº de fls.

    \printbibliography[heading=subbibliography,notcategory=fullcited]

    \label{chap:historia-vida-sertoesend}

\end{refsection}
